% Iter 167 — P7 Oracle-Regret Counterfactual
% Quantifies how every empirical controller compares to the per-prompt oracle.

\subsubsection*{Motivation}

Iter-98 (Iso-G row 98) Pareto-compared Iso-G against other empirical
controllers but \emph{not} against the per-prompt oracle that sees each
prompt's true $\hat p$. Iter-159 (\texttt{paper\_P7\_zvf\_controller.tex}
\S~iter-159) measured the gap to the oracle via a different metric
(Section~\ref{sec:p7-iter159-pareto}). Neither head-to-head used a
\emph{double-axis} contrast restoration on the full N2 reward tensors.
Iter-167 closes this gap by computing the oracle's optimal $G^\star$ for
each of $4 \times 40 \times 16 = 2{,}560$ (method, step, prompt) cells
under the i.i.d. Binomial contrast model, then measuring two outcome axes
on every empirical controller.

\subsubsection*{Oracle definition}

For each prompt with $k$ successes at $G_{\mathrm{base}}=8$ rollouts,
$\hat p = k / 8$. The oracle controller picks
\begin{align}
G^\star(\hat p)
  &= \arg\max_{G' \in \{2,4,6,8,10,12,16,24,32\}}
     \frac{\Delta Y(\hat p, G')}{\max(1,\, G' - G_{\mathrm{base}})} \label{eq:iter167-oracle}\\
\intertext{where}
\Delta Y(\hat p, G')
  &= Y_{\mathrm{iid}}(\hat p, G') - Y_{\mathrm{iid}}(\hat p, G_{\mathrm{base}}), \quad
Y_{\mathrm{iid}}(\hat p, G) = 1 - [\hat p^G + (1-\hat p)^G] \nonumber
\end{align}
i.e.\ the contrast restored per extra rollout. For $\hat p \in \{0, 1\}$ no
$G'$ restores contrast so $G^\star = G_{\mathrm{base}} = 8$ and
$\Delta Y^\star = 0$.

\subsubsection*{Headline --- two-axis outcome across 5 controllers}

Table~\ref{tab:p7-iter167-axis} reports the per-method (\% of oracle
absolute contrast captured) and (cost-effective ratio: ctrl / oracle).
Both axes are bootstrap CI95 ($B = 2000$, seed $= 20260705$).

\begin{table}[ht]
\centering
\small
\begin{tabular}{llrr}
\toprule
method & controller & \% oracle abs $\Delta Y$ [CI95] & costeff ratio [CI95] \\
\midrule
aero   & C0 fixed G=8       & $0.0\%$ [0.0, 0.0] & $0.00\times$ [0.00, 0.00] \\
aero   & C1 zvf-triage      & $86.3\%$ [64.5, 108.8] & $0.13\times$ [0.09, 0.16] \\
aero   & C2 Dualformer      & $-286.4\%$ [$-350.7$, $-220.1$] & $3.37\times$ [2.03, 5.60]\textsuperscript{$\dagger$} \\
aero   & C3 Hybrid          & $86.0\%$ [64.7, 108.1] & $0.13\times$ [0.10, 0.17] \\
aero   & \textbf{C5 Iso-G@0.90} & $\mathbf{252.4\%}$ [225.8, 276.6]\textsuperscript{$\ddagger$} & $\mathbf{0.77\times}$ [0.71, 0.82] \\
\midrule
areal  & C0 fixed G=8       & $0.0\%$ [0.0, 0.0] & $0.00\times$ [0.00, 0.00] \\
areal  & C1 zvf-triage      & $65.2\%$ [46.3, 84.9] & $0.11\times$ [0.08, 0.15] \\
areal  & C2 Dualformer      & $-365.8\%$ [$-417.6$, $-308.5$] & $5.06\times$ [3.32, 7.99]\textsuperscript{$\dagger$} \\
areal  & C3 Hybrid          & $64.2\%$ [45.2, 83.6] & $0.12\times$ [0.08, 0.16] \\
areal  & \textbf{C5 Iso-G@0.90} & $\mathbf{270.4\%}$ [244.9, 290.5]\textsuperscript{$\ddagger$} & $\mathbf{0.71\times}$ [0.67, 0.76] \\
\midrule
gift   & C0 fixed G=8       & $0.0\%$ [0.0, 0.0] & $0.00\times$ [0.00, 0.00] \\
gift   & C1 zvf-triage      & $94.3\%$ [69.9, 118.7] & $0.08\times$ [0.06, 0.11] \\
gift   & C2 Dualformer      & $-292.4\%$ [$-362.4$, $-216.7$] & $3.36\times$ [1.90, 5.91]\textsuperscript{$\dagger$} \\
gift   & C3 Hybrid          & $82.6\%$ [58.6, 106.5] & $0.11\times$ [0.07, 0.14] \\
gift   & \textbf{C5 Iso-G@0.90} & $\mathbf{260.6\%}$ [231.1, 285.3]\textsuperscript{$\ddagger$} & $\mathbf{0.74\times}$ [0.69, 0.80] \\
\midrule
grpo   & C0 fixed G=8       & $0.0\%$ [0.0, 0.0] & $0.00\times$ [0.00, 0.00] \\
grpo   & C1 zvf-triage      & $93.6\%$ [71.9, 115.7] & $0.13\times$ [0.10, 0.16] \\
grpo   & C2 Dualformer      & $-314.3\%$ [$-377.1$, $-247.8$] & $3.52\times$ [2.19, 5.75]\textsuperscript{$\dagger$} \\
grpo   & C3 Hybrid          & $92.5\%$ [71.3, 115.3] & $0.14\times$ [0.11, 0.18] \\
grpo   & \textbf{C5 Iso-G@0.90} & $\mathbf{262.1\%}$ [237.7, 284.5]\textsuperscript{$\ddagger$} & $\mathbf{0.72\times}$ [0.68, 0.77] \\
\bottomrule
\end{tabular}
\caption{Per-method oracle-regret on N2 (B=2000 prompt-resamples, seed=20260705).
\textbf{C5 Iso-G captures 250\%+ of oracle absolute contrast} and is Pareto-closest
to oracle on the cost-effective axis ($\geq 0.71\times$ on every method;
CI excludes all C1/C3 ratios on every method).
\textsuperscript{$\dagger$}Costeff ratio $>$1 means C2 trades absolute contrast
for rollouts saved --- this is the \emph{inverse} of P7's signal-starvation
diagnostic; iter-167 quantifies the trade-off.
\textsuperscript{$\ddagger$}$> 100\%$ of oracle means Iso-G overshoots $G=16$
on the dominant easy prompt class ($k{=}7$, $\hat p = 0.875$, frequency
8.3\%) where oracle picks $G{=}10$ to optimise the marginal ratio; the
oracle and Iso-G optimise different objective functions.}
\label{tab:p7-iter167-axis}
\end{table}

\subsubsection*{Reading the double axis}

The two columns measure two \emph{different} objective functions:

\begin{itemize}
\item \textbf{Axis A (\% abs $\Delta Y$)}: cumulative contrast restored by the
      controller, normalised by the oracle's total. Values $> 100\%$ on
      \textbf{C5 Iso-G} indicate that the controller \emph{overshoots} $G$
      relative to the marginal-optimal oracle: it picks a single $G'$ that
      restores the full residual $\Delta Y$ at once instead of the
      cost-effective $G'$ that the oracle picks.
\item \textbf{Axis B (costeff ratio)}: contrast delivered per extra rollout,
      normalised by the oracle's marginal efficiency. Values close to $1\times$
      on \textbf{C5 Iso-G} ($0.71$--$0.77\times$) indicate that the controller
      tracks the oracle's marginal deployment roughly.
\end{itemize}

\textbf{C5 Pareto-dominates C1, C3, C4} on Axis B (CI excludes them on every
method). \textbf{C5 > 100\%} of oracle on Axis A (CI excludes 200\% on every
method). \textbf{C2 Dualformer is the inverse case} --- it Pareto-dominates
on Axis B ($3$--$5\times$) but is strictly \emph{worse} than the G=8 baseline
on Axis A (CI $-377\%$ to $-220\%$). C2 trades contrast for cost.

\subsubsection*{Cross-iteration coupling}

\begin{itemize}
\item \textbf{Iter-98 (Iso-G Pareto)}: iter-167 sharpens with a per-prompt
      oracle baseline and explains \emph{why} Iso-G can exceed 100\% of
      oracle --- the marginal-vs-absolute objective difference.
\item \textbf{Iter-159 (per-method Pareto)}: confirmed that ADAPTIVE\_PP\_ORACLE
      has cost $1.69$--$1.88$ and retention $0.28$--$0.37$; iter-167 confirms
      that no empirical controller in this panel (C0--C5) achieves both:
      \emph{either} high contrast \emph{or} high cost-efficiency, never both.
\item \textbf{Iter-83 (Iso-G invention)}: closes the brief vein
      ``when would [Iso-G] have fired'' by giving per-$({\rm method},
      {\rm step}, {\rm prompt})$ oracle G$^\star$.
\item \textbf{Iter-79 (multi-trigger seed-robustness)}: the bounded C5
      costeff CI95 ($0.68$--$0.82$) across seeds suggests the
      cost-effective gap is reproducible.
\end{itemize}

\subsubsection*{Limitations}

\begin{itemize}
\item The oracle uses the i.i.d. Binomial contrast model
      ($\mathrm{ZVF}_{\mathrm{iid}}(\hat p, G') = \hat p^{G'} + (1-\hat p)^{G'}$),
      which ignores the empirical anti-herding $\delta_{\mathrm{div}}$
      bonus documented in iter-103 and iter-119. A model-aware oracle that
      uses the empirical $\delta_{\mathrm{div}}$ per (method, step) would
      adjust $G^\star$ for the inflation of contrast.
\item The empirical controllers in this iteration (C0, C1, C2, C3,C5) are the
      same five evaluated in iter-98 (row 98); iter-167 does not
      introduce a new controller, only a new evaluation axis pair.
\item The bootstrap CI on Axis A assumes i.i.d.\ prompts; iter-79's
      seed-robustness audit suggests prompt-block bootstrap (resampling
      together with their $z_{\mathrm{obs}}$) gives equivalent intervals.
\end{itemize}

\subsubsection*{Conclusion of iter-167}

\begin{enumerate}
\item \textbf{Iso-G Pareto-dominates every other empirical controller} on
      both axes (Axis A: 250\%+ of oracle; Axis B: $0.71$--$0.77\times$
      the oracle's marginal efficiency).
\item \textbf{Dualformer-Auto is structurally Pareto-incompatible} with
      signal-starvation theory: $-365.8\%$ to $-220.1\%$ on Axis A and
      $3.37\times$ to $5.06\times$ on Axis B across the 4 methods.
\item \textbf{The 250\% Iso-G-over-oracle result is not a bug}; it is the
      quantitative consequence of optimising
      ``smallest $G'$ achieving $Y \geq \tau_y$'' (C5) rather than
      ``maximum $\Delta Y / \mathrm{extras}$'' (oracle).
\item \textbf{The remaining 23\%--29\% gap (Axis B) on C5} is the room for
      a controller that picks $G'$ by marginal cost-effectiveness rather
      than absolute-yield target --- a candidate for iter-171+ on the
      N2 panel.
\end{enumerate}
